% =========================================================================
% English translation (generated by AI using Claude Opus 4.8 (Max effort)).
% This file was translated from Hungarian to English by AI.
% DISCLAIMER: Please review against the original Hungarian .tex files, before relying on it!
% SOURCE: 13. Számításelmélet.tex
% Note: Hungarian problem names were translated to their standard English
%       names (Elérhetőség->Reachability, Teljes részgráf->Clique,
%       Független csúcshalmaz->Independent set, Csúcslefedés->Vertex cover,
%       Utazóügynök->Travelling salesman, L_átló->L_diag, L_üres->L_empty).
% =========================================================================
\documentclass[margin=0px]{article}

\usepackage{listings}
\usepackage[utf8]{inputenc}
\usepackage{graphicx}
\usepackage{float}
\usepackage[a4paper, margin=0.7in]{geometry}
\usepackage{amsthm}
\usepackage{amssymb}
\usepackage{fancyhdr}
\usepackage{setspace}

\onehalfspacing

\makeatletter
\renewcommand\paragraph{%
	\@startsection{paragraph}{4}{0mm}%
	{-\baselineskip}%
	{.5\baselineskip}%
	{\normalfont\normalsize\bfseries}}
\makeatother

\renewcommand{\figurename}{Figure}

\newenvironment{tetel}[1]{\paragraph{#1 \\}}{}

\pagestyle{fancy}
\lhead{\it{CS BSc Final Exam Topics}}
\rhead{13. Theory of computation}

\title{\textbf{{\Large ELTE IK - Computer Science BSc} \vspace{0.2cm} \\ {\huge Final Exam Topics}} \vspace{0.3cm} \\ 13. Theory of computation}
\author{}
\date{}

\begin{document}
\maketitle

\begin{center}
\fbox{\parbox{0.92\textwidth}{\small \textit{Note: This document was translated from Hungarian to English by AI. Please verify the information against the original Hungarian source before relying on it.}}}
\end{center}

\begin{tetel}{Theory of computation}
    The Turing machine and the Church-Turing thesis. Variants of Turing machines: multitape, nondeterministic, counting, offline. Recursive and recursively enumerable languages. Undecidable problems. Time and space complexity classes: P, NP, PSPACE. NP-complete problems.
\end{tetel}

\section{Computability}

\subsubsection{Algorithm models}

\begin{itemize}
    \item	\textbf{Gödel}: recursive functions (primitive recursive functions in 1931, then more general in 1934)

    \item	\textbf{Church}: $\lambda$-calculus, $\lambda$-definable functions: equivalent to the recursive functions (proven)

    \item	\textbf{Turing}: Turing machine (1936); the $\lambda$-definable and the Turing-machine-computable functions coincide (proven)

\end{itemize}

\noindent \textbf{Church-Turing thesis}: The various mathematical models of computability all define the class of effectively
computable functions.

\subsubsection{Concepts}

By a computational problem we mean a question formulated in the language of mathematics
to which we want to give the answer with an algorithm. To almost every problem of
practical life, using an appropriate abstraction, a computational problem can be assigned.\\

\noindent A problem together with its associated concrete input is called an instance of the
problem.\\

A special computational problem is the decision problem. In this case the question related to the problem
is a yes/no question, so for an instance of the problem the answer will be ``yes'' or ``no''.\\

A computational problem can be represented by a function $f : A \to B$. The set $A$
contains the individual inputs of the problem, typically encoded in a word over an appropriate alphabet,
while the set $B$ contains the answers given to the inputs, likewise encoded in a word over some suitable alphabet. Naturally,
if it is a decision problem, then the range of $f$, that is, $B$,
is a two-element set: $\left\{yes, no\right\}$, $\left\{1, 0\right\}$, etc.\\

\noindent \textbf{Computable function}: A function $f : A \to B$ is called \textit{computable}
if for every element $x \in A$ the function value $f(x) \in B$ is computable by some
algorithmic model.\\

\noindent \textbf{Solvable, decidable problem}: A computational problem is \textit{solvable}
(in the case of a decision problem we say that it is \textit{decidable}) if the function determined by it
is computable.\\

\noindent \textbf{Time requirement of algorithms}: Let $f,g: \mathbb{N} \to \mathbb{N}$ be functions, where
$\mathbb{N}$ is the set of natural numbers. We say that $f$ grows at most as fast as $g$
(notation: $f(n) = \mathcal{O}(g(n))$) if $\exists c>0$ and $n_{0} \in \mathbb{N}$ such that
$f(n) \leq c * g(n) \ \forall n \geq n_{0}$. The notation $f(n) = \Omega(g(n))$ denotes that $g(n) = \mathcal{O}(f(n))$
holds, and $f(n) = \Theta(g(n))$ denotes that both $f(n) = \mathcal{O}(g(n))$ and $f(n) = \Omega(g(n))$ hold.\\

\noindent \textbf{Example}: $3n^{3} + 5n^{2} + 6 = \mathcal{O}(n^{3})$, $n^{k} = \mathcal{O}(2^{n}) \ \forall k \geq 0$, etc.\\

\noindent \textbf{Theorem}: Every polynomial function grows more slowly than any exponential function,
that is, for every polynomial $p(n)$ and every $c>0$ there is an integer $n_{0}$ such that for $\forall n \geq n_{0}$, $p(n) \leq 2^{cn}$\\

\noindent \textbf{Correspondence of a computational problem to a decision problem}: Consider a computational problem $P$
and let $f: A \to B$ be the function determined by $P$. Then a decision problem $P'$ can be given for $P$ such that
$P'$ is decidable exactly when $P$ is computable. Namely, let us pair, for every element $a \in A$, the elements $a$ and $f(a)$,
and encode the pairs thus obtained each in a word. After this let $P'$ be the formal language formed from the words thus obtained.
It is obvious that if for every element $a \in A$ and $b \in B$ the containment $(a,b) \in P'$ is decidable (that is, $P'$ is decidable), then
$P$ is computable and vice versa.
Because of this correspondence, in what follows we typically deal with decision problems.

\section{Turing machines}

Similarly to the finite automaton or the pushdown automaton, the Turing machine is also
a device with finitely many states. The Turing machine works on a two-way
infinite tape. The tape is divided into cells; this is essentially the (unlimited) memory of the machine.
Initially on the tape there is only the input word, one letter on each cell. The other cells of the tape are
filled with so-called blank or space ($\sqcup$) symbols. Initially the so-called read-write
head of the machine stands on the first letter of the input word and the machine is in its initial state.
The machine can move the read-write head arbitrarily on the tape. It is furthermore able to
read out and rewrite the content of the tape at the head position. The machine has two
distinguished states, the states $q_{i}$ and $q_{n}$. If it gets into these states,
then it accepts respectively rejects the input word.
Formally we define the Turing machine in the following way.\\

\noindent \textbf{The formal definition of the Turing machine}: The Turing machine is a system
$M = (Q, \Sigma, \Gamma, \delta, q_{0}, q_{i}, q_{n})$, where:

\begin{itemize}

    \item	$Q$ is the finite, nonempty set of states,

    \item	$q_{0}, q_{i}, q_{n} \in Q$, $q_{0}$ is the initial state, $q_{i}$ the accepting state,
          and $q_{n}$ the rejecting state,

    \item	$\Sigma$ and $\Gamma$ are alphabets, the alphabet of the input symbols and of the tape symbols such that
          $\Sigma \subseteq \Gamma$ and $\Gamma - \Sigma$ contains a special $\sqcup$ symbol,

    \item	$\delta : (Q - \left\{q_{i},q_{n}\right\}) \times \Gamma \to Q \times \Gamma \times \left\{L, R, S\right\}$
          is the transition function.

\end{itemize}

As in the case of pushdown automata, the phases of operation of a Turing machine $M$ can also be described with configurations.\\

\noindent \textbf{Configuration of a Turing machine}: A configuration of the Turing machine $M$ is a word $uqv$ where
$q \in Q$ and $u, v \in \Gamma^{*}$, $v \not = \varepsilon$. This configuration reflects that state of $M$
when the content of the tape is $uv$ (before and after $uv$ there is only $\sqcup$ on the tape), the machine is in state $q$,
and the read-write head points to the first letter of $v$. We denote the set of all configurations of $M$ by $\mathcal{C}_{M}$.\\

\noindent \textbf{Initial configuration of a Turing machine}: The initial configuration of $M$ is a word $q_{0}u\sqcup$ where
$u$ contains only letters from $\Sigma$.\\

\noindent \textbf{Configuration transition of a Turing machine}: The configuration transition of $M$ is a relation
$\vdash \subseteq \mathcal{C}_{M} \times \mathcal{C}_{M}$, which we define as follows.
Let $uqav$ be a configuration where $a \in \Gamma$ and $u, v \in \Gamma^{*}$. We distinguish the following three
cases:

\begin{enumerate}

    \item	If $\delta(q,a) = (r, b, S)$, then $uqav \vdash urbv$.

    \item	If $\delta(q,a) = (r, b, R)$, then $uqav \vdash ubrv'$, where $v' = v$ if $v \not = \varepsilon$, otherwise
          $v' = \sqcup$.

    \item	If $\delta(q,a) = (r, b, L)$, then $uqav \vdash u'rcbv$, where $u'c = u$ for some $u' \in \Gamma^{*}$
          and $c \in \Gamma$, if $u \not = \varepsilon$, otherwise $u' = \varepsilon$, $c = \sqcup$.

\end{enumerate}

We say that $M$ reaches the configuration $C'$ from the configuration $C$ in finitely many steps (notation $C \vdash^{*} C'$)
if there exist $n \geq 0$ and a configuration sequence $C_{1}, ... C_{n}$ such that $C_{1} = C$, $C_{n} = C'$ and for every
$1 \leq i < n$, $C_{i} \vdash C_{i+1}$.

If $q \in \left\{q_{i}, q_{n}\right\}$, then we say that the configuration $uqv$ is a halting
configuration. Furthermore, in the case $q = q_{i}$ we speak of an accepting, and in the case $q = q_{n}$ of a rejecting
configuration.\\

\noindent \textbf{Language recognized by a Turing machine}: The language recognized by the Turing machine $M$ (notation $L(M)$)
is the set of those words $u \in \Sigma^{*}$ for which $q_{0}u\sqcup \vdash^{*} xq_{i}y$ holds
for some words $x,y \in \Gamma^{*}$, $y \not = \varepsilon$.

\begin{figure}[H]
    \centering
    \includegraphics[width=0.6\linewidth]{../../13. Számításelmélet (és Logika)/img/turinggep_pelda}
    \caption{A Turing machine recognizing $L = \left\{u\#u \ | \ u \in \left\{0,1\right\}^{+}\right\}$.}
    \label{fig:turinggep_pelda}
\end{figure}

\noindent \textbf{Equivalence of Turing machines}: Two Turing machines are called equivalent if they recognize the same language.\\

\noindent \textbf{Turing-recognizable language, the class of recursively enumerable languages}:
A language $L \subseteq \Sigma^{*}$ is Turing-recognizable if $L = L(M)$ for some Turing machine $M$.
Turing-recognizable languages are also usually called \textit{recursively enumerable}.
The class of recursively enumerable languages is denoted by $RE$.\\

\noindent \textbf{Turing-decidable language, the class of recursive languages}:
A language $L \subseteq \Sigma^{*}$ is Turing-decidable if there exists a Turing machine
that on every input reaches a halting configuration and recognizes $L$.
Turing-decidable languages are also usually called \textit{recursive}.
The class of recursive languages is denoted by $R$.\\

\noindent \textbf{Running time, time requirement of a Turing machine}:
Consider a Turing machine $M = (Q, \Sigma, \Gamma, \delta, q_{0}, q_{i}, q_{n})$ and one of its
input words $u \in \Sigma^{*}$. We say that the running time (time requirement) of $M$ on the word $u$ is $n$ ($n \geq 0$),
if $M$ can reach a halting configuration in $n$ steps from the initial configuration $q_{0}u\sqcup$. If there is no such
number, then the running time of $M$ on the word $u$ is infinite.

Let $f : \mathbb{N} \to \mathbb{N}$ be a function. We say that the time requirement of $M$ is $f(n)$ (or that $M$
is an $f(n)$ time-bounded machine) if for every input word $u \in \Sigma^{*}$ the time requirement of $M$ on the word $u$ is at most $f(l(u))$.

\subsubsection{Multitape Turing machines}

Multitape Turing machines, naturally, have more than one tape.
To each tape belongs a read-write head, which can move on the tape independently of each other.\\

\noindent \textbf{Definition of a multitape Turing machine}: Let $k > 1$. A $k$-tape Turing machine is a system
$M = (Q, \Sigma, \Gamma, \delta, q_{0}, q_{i}, q_{n})$, where the components except for $\delta$
coincide with the components of the single-tape Turing machine, and $\delta$ is given as follows.
$\delta : (Q - \left\{q_{i},q_{n}\right\}) \times \Gamma^{k} \to Q \times \Gamma^{k} \times \left\{L, R, S\right\}^{k}$.
Let $q, p \in Q$, $a_{1}, a_{2}, ... , a_{k}, b_{1}, b_{2}, ..., b_{k} \in \Gamma$ and
$D_{1}, D_{2}, ..., D_{k} \in \left\{L, R, S\right\}$. If $\delta(q,a_{1}, a_{2}, ... , a_{k}) = (p,b_{1}, b_{2}, ..., b_{k}, D_{1}, D_{2}, ..., D_{k})$, then the machine, from state $q$, if on its tapes it reads respectively the letters
$a_{1}, a_{2}, ... , a_{k}$, can transition into state $p$, while it rewrites the letters
$a_{1}, a_{2}, ... , a_{k}$ to the letters $b_{1}, b_{2}, ... , b_{k}$ and moves the heads on the tapes in the
$D_{1}, D_{2}, ... , D_{k}$ directions.

The definition of the configuration of the multitape Turing machine, the configuration transition, as well as the
recognized or decided language is the natural generalization of the single-tape case.
We also define the time requirement of the multitape Turing machine similarly to the single-tape one.\\

\noindent \textbf{Equivalence of multitape and single-tape machines}: For every $k$-tape, $f(n)$ time-bounded Turing machine
there is an equivalent $\mathcal{O}(n*f(n))$ time-bounded single-tape Turing machine.

\subsubsection{Nondeterministic Turing machines}

The state function of a nondeterministic Turing machine $M$ is of the form
$\delta : (Q - \left\{q_{i},q_{n}\right\}) \times \mathcal{P}(\Gamma \to Q \times \Gamma \times \left\{L, R\right\})$.
So $M$ can transition from each of its configurations into some (possibly zero) different configurations.
In this way the computation sequences of $M$ on a word $u$ can be represented by a tree. The root of the tree is the initial configuration of $M$,
and its vertices are the configurations of $M$. Every leaf of the tree corresponds to a computation sequence of $M$ on
$u$. $M$ accepts $u$ if some leaf of the tree is an accepting configuration.
Let us call the tree just described the nondeterministic computation tree of $M$ on
$u$. The language recognized by $M$ can be defined similarly to the deterministic case,
and the language decided by the machine as follows.\\

\noindent \textbf{Language decided by a nondeterministic Turing machine}: We say that a nondeterministic
Turing machine $M$ decides a language $L \subseteq \Gamma^{*}$ if it recognizes it, and for every word $u \in \Sigma^{*}$
the computation sequences of $M$ are finite and lead to an accepting or rejecting configuration.\\

\noindent \textbf{Time requirement of a nondeterministic Turing machine}: Let $f : \mathbb{N} \to \mathbb{N}$ be a function, $M$
a nondeterministic Turing machine. The time requirement of $M$ is $f(n)$ if on an input $u$ of length $n$ there are no
computation sequences of $M$ longer than $f(n)$, that is, the computation tree of $M$ on $u$ is at most $f(n)$ high.\\

\noindent \textbf{Equivalence of deterministic and nondeterministic Turing machines}:
For every nondeterministic Turing machine $M$ an equivalent deterministic Turing machine $M'$ can be given.
Furthermore, if $M$ is of time requirement $f(n)$ for some function $f : \mathbb{N} \to \mathbb{N}$, then
$M'$ is of time requirement $2^{\mathcal{O}(f(n))}$.


\section{Undecidable problems}

In this chapter we show that although the Turing machine is the most general possible
algorithm model, there are nevertheless problems that cannot be computed with a Turing machine.\\

\noindent \textbf{Reminder}: The class of recursively enumerable (Turing-recognizable) languages is denoted by $RE$,
the class of recursive (Turing-decidable) languages by $R$.\\

It is clear that $R \subseteq RE $. Our goal is to show that $R$ is a proper subset of $RE$,
that is, there is a language (problem) that is Turing-recognizable but not decidable.

We will examine only those Turing machines whose input alphabet is the set $\left\{0, 1\right\}$.
This does not mean a restriction of generality, since if we find a language over $\left\{0, 1\right\}$
that cannot be decided with such a Turing machine, then this language cannot be decided at all.\\

\subsubsection{Encoding of Turing machines}

The words over $\left\{0, 1\right\}$ can be enumerated (that is, they are countable). Indeed, consider the enumeration
in which the words follow each other by their length, and of two words of equal length the one that
precedes the other in alphabetical order comes first. In this way
an enumeration of the elements of the set $\left\{0, 1\right\}^{*}$ takes the following shape: $w_{1} = \varepsilon$,
$w_{2} = 0$, $w_{3} = 1$, $w_{4} = 00$, $w_{5} = 01$ and so on. In this chapter, therefore, by the word
$w_{i}$ we denote the $i$-th element of $\left\{0, 1\right\}^{*}$.

Let furthermore $M$ be a Turing machine over the input alphabet $\left\{0, 1\right\}$. There is a number $k > 0$ such that
$Q$ can be written in the form $Q = \left\{p_{1}, ... p_{k}\right\}$, where $p_{1} = q_{0}$, $p_{k-1} = q_{i}$, $p_{k} = q_{n}$.
Furthermore, there is a number $m > 0$ such that $\Gamma$ can be written in the form $\Gamma = \left\{X_{1}, ... X_{m}\right\}$,
where $X_{1} = 0$, $X_{2} = 1$, $X_{3} = \sqcup$, and $X_{4}, ... X_{m}$ are the further tape symbols of $M$.
Let us finally call the symbols $L, R, S$ (which denote directions) respectively $D_{1}$, $D_{2}$ and $D_{3}$.
After this a transition $\delta(p_{i},X_{j}) = (p_{r}, X_{s}, D_{t})$ ($0 \leq i,r \leq k$, $1 \leq j,s \leq m$ and $1 \leq t \leq 3$)
of $M$ can be encoded with the word $0^{i}10^{j}10^{r}10^{s}10^{t}$. Since the length of every $0$-block is at least $1$, in the word
encoding the transition the subword $11$ does not appear. So the words encoding all transitions of M can be concatenated into a word
in which the transitions are separated from each other by the subword $11$. The word thus obtained encodes $M$ itself.\\

In what follows we denote by $M_{i}$ the Turing machine encoded by the word $w_{i}$ ($i \geq 1$). If
$w_{i}$ is not the encoding of a Turing machine described above, then let us regard $M_{i}$ as one that, for every input,
immediately goes into the state $q_{n}$, that is, $L(M_{i}) = \emptyset$.\\

Later on we will need to encode a pair of a Turing machine and input $(M, w)$ in a word
over $\left\{0, 1\right\}$. Since the encoding of Turing machines cannot contain
$111$, therefore the encoding of $(M, w)$ is the following: after the encoding of $M$ we write $111$, then after it $w$.

\subsubsection{A non-recursively-enumerable language}

\noindent \textbf{The $L_{diag}$ language}: The $L_{diag}$ language contains the binary encodings of those Turing machines
over $\left\{0, 1\right\}$ that do not accept their own encoding as the input word, that is,
$L_{diag} = \left\{w_{i} \ | \ i \geq 1, w_{i} \notin L(M_{i}) \right\}$\\

\noindent \textbf{Theorem}: $L_{diag} \notin RE$.

\subsubsection{A recursively enumerable but undecidable language}

\noindent \textbf{The $L_{u}$ language}: Consider the set of those pairs $(M, w)$ (encoded in an appropriate binary word),
where $M$ is a Turing machine over the input alphabet $\left\{0, 1\right\}$, and $w$ is a word over $\left\{0, 1\right\}$
such that $w \in L(M)$, that is, $M$ accepts $w$. We denote this language by $L_{u}$.
$L_{u} = \left\{\langle w_{i},w_{j} \rangle \ | \ i, j \geq 1, w_{j} \in L(M_{i}) \right\}$\\

\noindent \textbf{Theorem}: $L_{u} \in RE$.\\

\noindent \textbf{Theorem}: $L_{u} \notin R$.

\subsubsection{Further theorems}

\begin{enumerate}
    \item	Let $L$ be a language. If $L, \bar{L} \in RE$, then $L \in R$. Consequence: the recursively enumerable
          languages are not closed under complementation.

    \item	If $L \in R$, then $\bar{L} \in R$, that is, the recursive languages are closed under complementation.
\end{enumerate}

\subsubsection{Further undecidable problems}

\noindent \textbf{Computable function}: Let $\Sigma$ and $\Delta$ be two alphabets and $f$ a function mapping from $\Sigma^{*}$
into $\Delta^{*}$. We say that $f$ is computable if there is a Turing machine $M$ such that,
starting $M$ with a word $w \in \Sigma^{*}$ on its input, $M$ halts so that on its tape there is the
word $f(w) \in \Delta^{*}$.\\

\noindent \textbf{Reduction of decision problems}: Let $L_{1} \subseteq \Sigma^{*}$ and $L_{2} \subseteq \Delta^{*}$
be two decision problems. $L_{1}$ is reducible to $L_{2}$ ($L_{1} \leq L_{2}$) if there is a
computable function $f : \Sigma^{*} \to \Delta^{*}$ such that for every word $w \in \Sigma^{*}$,
$w \in L_{1}$ holds exactly when $f(w) \in L_{2}$ also holds.\\

\noindent \textbf{Theorem}: Let $L_{1} \subseteq \Sigma^{*}$ and $L_{2} \subseteq \Delta^{*}$
be two decision problems and suppose that $L_{1}$ is reducible to $L_{2}$. Then the following statements are true:

\begin{enumerate}
    \item	If $L_{1}$ is undecidable, then so is $L_{2}$.

    \item	If $L_{1} \notin RE$, then $L_{2} \notin RE$.
\end{enumerate}

\noindent \textbf{The halting problem}: Let $L_{h} = \left\{\langle M, w \rangle \ |\ M \ halts \ on \ the \ input \ w \right\}$,
that is, $L_{h}$ contains, encoded, those Turing-machine-and-input pairs $\langle M, w \rangle$ for which $M$ halts on the input $w$.
$L_{h}$ is undecidable ($L_{u}$ is reducible to $L_{h}$), but $L_{h} \in RE$.\\

\noindent \textbf{The $L_{empty}$ problem}: Let $L_{empty} = \left\{\langle M \rangle \ |\ L(M) = \emptyset \right\}$.
$L_{empty}$ is undecidable ($L_{u}$ is reducible to $L_{empty}$), and $L_{empty} \notin RE$.\\

\noindent \textbf{(Non-trivial) property of recursively enumerable languages}: If $\mathcal{P}$ is a set of recursively
enumerable languages, then $\mathcal{P}$ is a property of the recursively enumerable languages. If $\mathcal{P} \not = \emptyset$ and
$\mathcal{P} \not = RE$, then $\mathcal{P}$ is a non-trivial property of the recursively enumerable languages.\\

\noindent \textbf{Rice's theorem}:
For a given property $\mathcal{P}$ let us denote by $L_{\mathcal{P}}$ the set of encodings of those Turing machines that
recognize a language in $\mathcal{P}$. If $\mathcal{P}$ is a non-trivial property of the recursively enumerable languages, then
$L_{\mathcal{P}}$ is undecidable.\\

\noindent \textbf{Post Correspondence Problem (briefly PCP)}: We define the PCP problem as follows. Let
$\Sigma$ be an alphabet containing at least two letters and let $D = \left\{[\frac{u_{1}}{v_{1}}], ..., [\frac{u_{n}}{v_{n}}]\right\}$
be a set of dominoes, in which $n \geq 1$ and $u_{1}, ..., u_{n}, v_{1}, ..., v_{n} \in \Sigma^{+}$. The question is whether there is an
index sequence $1 \leq i_{1}, ..., i_{m} \leq m$ ($m \geq 1$) for which it holds that, writing the dominoes
$[\frac{u_{i_{1}}}{v_{i_{1}}}], ..., [\frac{u_{i_{m}}}{v_{i_{m}}}]$ next to each other, the same word is obtained below and above, that is,
$u_{i_{1}} ... u_{i_{m}} = v_{i_{1}} ... v_{i_{m}}$. In this case the above domino sequence is called a solution of D.

As a formal language we can define the PCP as follows:
PCP $= \left\{ \langle D \rangle \ |\ D \ has \ a \ solution \right\}$. PCP is undecidable.

\section{Complexity theory}

The goal of complexity theory is the classification of solvable (and within these, decidable) problems according to the
amount of resources (typically time and space) needed for the solution.

\subsubsection{Time complexity concepts}

\noindent \textbf{TIME}: Let $f : \mathbb{N} \to \mathbb{N}$ be a function. \textbf{TIME}($f(n)$) $= \left\{L \ | \ L \ is \ decidable \ by \ an \ \mathcal{O}(f(n)) \ time \ Turing \ machine \right\}$\\

\noindent \textbf{P} $=\bigcup_{k \geq 1}$ \textbf{TIME}($n^{k}$). So \textbf{P} contains those languages
that are decidable by a polynomial-time-bounded deterministic Turing machine. Such is, for example, the well-known \textsc{Reachability}
problem, whose input is a graph $G$ and its two distinguished vertices ($s$ and $t$). The question is whether there is a
path in $G$ from $s$ to $t$. If we regard the \textsc{Reachability} problem as a language, then we can write that\\

\textsc{Reachability} $= \left\{\langle G, s, t \rangle \ | \ there \ is \ a \ path \ in \ G \ from \ s \ to \ t\right\}$.\\

A deterministic Turing machine deciding the \textsc{Reachability} problem in polynomial time can easily be given, so
\textsc{Reachability} $\in$ \textbf{P}.\\

\noindent \textbf{NTIME}: Let $f : \mathbb{N} \to \mathbb{N}$ be a function.\\ \textbf{NTIME}($f(n)$) $= \left\{L \ | \ L \ is \ decidable \ by \ an \ \mathcal{O}(f(n)) \ time \ nondeterministic \ Turing \ machine \right\}$\\

\noindent \textbf{NP} $=\bigcup_{k \geq 1}$ \textbf{NTIME}($n^{k}$).
The problems in \textbf{NP} have a common property in the following sense.
If we consider an instance of a problem in \textbf{NP} and a possible ``proof''
that this instance is a ``yes'' instance of the given problem,
then the checking of the correctness of this proof can be done in polynomial time.
Accordingly, a nondeterministic Turing machine deciding a problem in \textbf{NP}
generally works so that it ``guesses'' a possible solution of the problem's input,
and checks in polynomial time whether the solution is correct.

Consider the \textsc{Sat} problem, which we define as follows. Given a propositional-logic CNF $\phi$.
The question is whether it is satisfiable. The proof that $\phi$ is satisfiable is a variable assignment
under which, evaluating $\phi$, we get a true value. An arbitrary variable assignment is therefore a possible proof of the
satisfiability of $\phi$. And the checking of whether this
assignment really makes $\phi$ true can be done in polynomial time. \textsc{Sat}
is a problem in \textbf{NP}.\\

\noindent It follows from the definitions that the containment \textbf{P} $\subseteq$ \textbf{NP} holds.

\subsubsection{NP-complete problems}

\noindent \textbf{Polynomial-time computable function}: Let $\Sigma$ and $\Delta$ be two alphabets and $f$ a function mapping from $\Sigma^{*}$
into $\Delta^{*}$. We say that $f$ is polynomial-time computable if it is computable by a polynomial-time
Turing machine.\\

\noindent \textbf{Polynomial-time reduction of decision problems}:
Let $L_{1} \subseteq \Sigma^{*}$ and $L_{2} \subseteq \Delta^{*}$ be two decision problems.
$L_{1}$ is polynomial-time reducible to $L_{2}$ ($L_{1} \leq_{p} L_{2}$) if $L_{1} \leq L_{2}$ and
the function $f$ used in the reduction is polynomial-time computable.\\

\noindent \textbf{Theorem}: Let $L_{1}$ and $L_{2}$ be two problems such that $L_{1} \leq_{p} L_{2}$. If $L_{2}$ is

\begin{enumerate}
    \item	in \textbf{P}, then $L_{1}$ is also in \textbf{P}.

    \item	in \textbf{NP}, then $L_{1}$ is also in \textbf{NP}.
\end{enumerate}

\noindent \textbf{NP-complete problem}: Let $L$ be a problem. We say that $L$ is \textbf{NP}-complete if

\begin{enumerate}
    \item	it is in \textbf{NP}, and

    \item	every further problem in \textbf{NP} is polynomial-time reducible to $L$.

\end{enumerate}

\noindent \textbf{Theorem}: Let $L$ be an \textbf{NP}-complete problem. If $L \in$ \textbf{P}, then \textbf{P} $=$ \textbf{NP}.\\

\noindent \textbf{Remark}: Currently we do \textbf{NOT} know of any \textbf{NP}-complete problem in \textbf{P}!!!\\

\noindent \textbf{Theorem}: Let $L_{1}$ be an \textbf{NP}-complete, and $L_{2}$ a problem in \textbf{NP}.
If $L_{1} \leq_{p} L_{2}$, then $L_{2}$ is also \textbf{NP}-complete.\\

\noindent \textbf{Cook's theorem}: \textsc{Sat} is \textbf{NP}-complete.\\

\noindent Let $k \geq 1$.
k\textsc{Sat} $= \left\{\langle \phi \rangle \ | \ in \ every \ clause \ of \ \phi \ there \ are \ k \ literals. \right\}$\\

\noindent \textbf{Theorem}: 3\textsc{Sat} is \textbf{NP}-complete, since \textsc{Sat} $\leq_{p}$ 3\textsc{Sat}.\\

\noindent \textsc{Clique} $= \left\{\langle G, k \rangle \ | \ G \ is \ a \ finite \ graph, k \geq 1, \ G \ has \ a \ subgraph \ with \ k
    \ vertices  \right\}$.
So \textsc{Clique} contains those pairs $G$ and $k$, encoded in words over an appropriate alphabet,
for which it holds that in $G$ there is a complete subgraph with $k$ vertices,
that is, a subgraph in which there is an edge between any two vertices.\\

\noindent \textsc{Clique} $= \left\{\langle G, k \rangle \ | \ G \ is \ a \ finite \ graph, k \geq 1, \ G \ has \ a \ subgraph \ with \ k
    \ vertices  \right\}$.
So \textsc{Clique} contains those pairs $G$ and $k$, encoded in words over an appropriate alphabet,
for which it holds that in $G$ there is a complete subgraph with $k$ vertices,
that is, a subgraph in which there is an edge between any two vertices.\\

\noindent \textsc{Independent set} $= \left\{\langle G, k \rangle \ | \ G \ is \ a \ finite \ graph, k \geq 1, \ G \ has \ an \ independent \ vertex \ set \ of \ k \ elements  \right\}$.
That is, \textsc{Independent set} contains those pairs $G$ and $k$ for which
it holds that in $G$ there are $k$ vertices, none of which is connected to
another.\\

\noindent \textsc{Vertex cover} $=$\\
$\left\{\langle G, k \rangle \ | \begin{array}{lr}
        \ G \ is \ a \ finite \ graph, k \geq 1, \ G \ has \ a \ vertex \ set \ of \ k \ elements, \\
        which \ contains \ at \ least \ 1 \ endpoint \ of \ every \ edge \ of \ G.
    \end{array}
    \right\}$.\\


\noindent \textsc{Clique}, \textsc{Independent set} and \textsc{Vertex cover} are \textbf{NP}-complete
(\textsc{Clique} $\leq_{p}$ \textsc{Independent set} $\leq_{p}$ \textsc{Vertex cover}).\\

\noindent \textsc{Travelling salesman} $=$\\
$\left\{\langle G, k \rangle \ |
    \begin{array}{lr}
        \ G \ is \ a \ finite \ undirected \ graph, \ on \ the \ edges \ a \
        positive \ integer \ weight \ each, \ and \\
        there \ is \ in \ G \ a \ Hamiltonian \ cycle \ of \ total \ weight \ at \ most \ k
    \end{array}
    \right\}$.\\

\noindent \textbf{Theorem}: The \textsc{Travelling salesman} problem is \textbf{NP}-complete.

\subsubsection{Space complexity}
We examine space complexity on a special, so-called offline Turing machine.\\

\noindent \textbf{Off-line Turing machine}: By an offline Turing machine we mean a multitape Turing machine
that can only read the tape containing the input, and can also write on the other, so-called work tapes.
Into the space requirement of the offline Turing machine only the area used on the work tapes is counted.\\

In what follows, by Turing machine we always mean an offline Turing machine. Now let us
define the language classes related to space complexity.\\

\noindent \textbf{SPACE}(f(n)) = $\left\{L|L \ is \ decidable \ by \ an \ \mathcal{O}(f(n)) \ space \
    deterministic \ Turing \ machine \right\}$\\

\noindent \textbf{NSPACE}(f(n)) = $\left\{L|L \ is \ decidable \ by \ an \ \mathcal{O}(f(n)) \ space \
    nondeterministic \ Turing \ machine \right\}$\\

\noindent \textbf{PSPACE} $=\bigcup_{k > 0}$ \textbf{SPACE}($n^{k}$)\\

\noindent \textbf{NPSPACE} $=\bigcup_{k > 0}$ \textbf{NSPACE}($n^{k}$)\\

\noindent \textbf{L} $=$ \textbf{SPACE}($\log_{2} n$) \\

\noindent \textbf{NL} $=$ \textbf{NSPACE}($\log_{2} n$) \\

\noindent \textbf{Savitch's theorem}: If $f(n) \geq \log n$, then \textbf{NSPACE}($f(n)$) $\subseteq$ \textbf{SPACE}($f^{2}(n)$).
\end{document}
